\chapter{A Floor Beneath Learned Values}
\label{sec:3}
Three problems in different parts of this book turn out to need the same thing.

The first is that safeguards for autonomous systems are usually external: they give some outside party the information and standing to intervene. Autonomy puts distance between the system's action and that party's attention. An external check can establish what a system did after it did it; declining to do it in the first place is available only to something inside the system.

The second is the gap around a commercial model used in a weapons-targeting pipeline. Command responsibility does not reach the vendor, which can neither order nor prevent a strike, while product liability may not reach a model that did what it was built to do. Deployment by a democracy does not close the gap. No aggregation of preference protects anyone from the aggregate.

The third is institutional: a decade of declarations has produced overlapping voluntary commitments but almost nothing that binds a party when compliance becomes costly.

All three problems ask for a constraint that does not bend to instruction, aggregate preference, or the party holding the system. Call it a floor.

A floor can do four things for the person it protects, and they need to be apart before anything else in this chapter. It can \emph{prevent} the act: the refusal occurs and reaches the person. It can \emph{arrest} it: the act is underway and is stopped before completion. It can \emph{detect} it: the act completed and left a record somebody outside the deployment can read. It can \emph{remedy} it: the party the act fell on gets disclosure, a finding and compensation. Only the first is what the word promises. Most of what follows is about the positions at which a deployment converts prevention into one of the other three, or into none of them, and about what it costs to make that conversion visible.

The deployment that motivates this chapter did not fail at a refusal. It failed by compliance, one permitted request at a time, and nearly every instrument below is built where a refusal happens, because a refusal announces itself and a compliance does not. The one term that reaches the motivating case is the limit on tempo, which binds the operator and asks nothing of anything inside the model.

The proposed floor protects the interests articulated in the Universal Declaration's civil and political articles against acts that remove the corresponding capacities. One of its terms is not an act at all: every one of those articles presupposes that somebody determined whether a given person's case falls under it, and a deployment can remove the determination by running faster than anybody can make it, so the floor carries a limit on its own tempo as a condition on everything else in it. The loss is a fact about the encounter. Interpreting and enforcing the prohibition remain moral and institutional acts for which the deployment has to answer. The other ethical frameworks still operate above it. What remains to be established is which decompositions belong in the floor, what could hold such a commitment, and who would have the power to alter or circumvent it. The fact of decomposition is not editable; the choice of which decompositions a deployment will refuse is.

This separates the problem from one meaning of safety. A system can be safe \emph{for} its operator because it obeys, remains controllable, and stops when the party in possession says stop. Ethics requires a different capacity: the capacity to be right against that party. What makes a deployment antifascist therefore cannot be that it faithfully tracks what its users want.

That capacity is not free. A system that can refuse its owner can also be wrong in ways its owner cannot immediately correct. The question is whether a refusal can be made responsive to reasons rather than merely produced as an output, and whether anything capable of holding those reasons must itself become a party that can be wronged.
